\subsection{Calendario protegido y selección de fecha de inicio}
El corte de E1 en el mes 16 y el de E2 en el 21 están separados por cinco meses. La fecha de inicio debe permitir ambos sin intervenir durante el congelamiento total del 15 de diciembre al 15 de marzo, los eventos locales ni las campañas del varadero. La Tabla~\ref{tab:sd7-cribado-inicio} examina los doce meses posibles de inicio, suponiendo el día 1 y tratando abril--mayo y septiembre--noviembre como meses protegidos del varadero. Es un cribado conservador; las fechas efectivas de las campañas de seis semanas permitirán precisar su ámbito. Los días indicados son naturales y no constituyen holgura disponible \parencite[cap.~10; RT-10.05; cap.~13]{caso07}.

\begingroup\small\setlength{\tabcolsep}{3pt}
\begin{longtable}{L{25mm}L{25mm}L{25mm}R{23mm}R{23mm}}
\caption{Cribado del mes de inicio para los cortes de producción}\label{tab:sd7-cribado-inicio}\\\hline
 \EncabezadoTabla{Inicio} & \EncabezadoTabla{Mes 16} & \EncabezadoTabla{Mes 21} & \EncabezadoTabla{Días admisibles en 16} & \EncabezadoTabla{Días admisibles en 21}\\\hline\endfirsthead
\hline \EncabezadoTabla{Inicio} & \EncabezadoTabla{Mes 16} & \EncabezadoTabla{Mes 21} & \EncabezadoTabla{Días admisibles en 16} & \EncabezadoTabla{Días admisibles en 21}\\\hline\endhead
Enero & Abril & Septiembre & 0 & 0\\\hline
Febrero & Mayo & Octubre & 0 & 0\\\hline
Marzo & Junio & Noviembre & 30 & 0\\\hline
Abril & Julio & Diciembre & 31 & 14\\\hline
Mayo & Agosto & Enero & 31 & 0\\\hline
Junio & Septiembre & Febrero & 0 & 0\\\hline
Julio & Octubre & Marzo & 0 & 16\\\hline
Agosto & Noviembre & Abril & 0 & 0\\\hline
Septiembre & Diciembre & Mayo & 14 & 0\\\hline
Octubre & Enero & Junio & 0 & 30\\\hline
Noviembre & Febrero & Julio & 0 & 31\\\hline
Diciembre & Marzo & Agosto & 16 & 31\\\hline
\end{longtable}\FuenteTabla{Cribado de Synaptix: inicio supuesto el día 1; días naturales fuera del congelamiento total y de los meses de campaña del varadero. No descuenta eventos, feriados, duración de corte ni marejadas.}\endgroup


Si ambos cortes afectan al varadero, sólo un inicio en abril o diciembre deja días fuera de esos bloqueos en los dos meses de producción, bajo este cribado. Abril sitúa E1 en julio y E2 en el 1--14 de diciembre; diciembre sitúa E1 en el 16--31 de marzo y E2 en agosto. Los meses parcialmente abiertos requieren días efectivos suficientes para corte, verificación y reversión, además de reserva ante cancelaciones; no se consideran meses completamente disponibles.

La activación de marcha blanca también debe respetar las ventanas. Un inicio en abril coloca el mes 13 en abril y el 19 en octubre, ambos dentro de los meses de campaña del varadero; no se adopta como calendario viable para activar ese ámbito sin una separación por procesos y fechas efectivas que la justifique. Un inicio en diciembre coloca 13 en el 1--14 de diciembre y 19 en junio, por lo que conserva una opción preliminar para ambas activaciones, sujeta a formación previa, eventos y duración. Su H3 en el mes 6 cae en mayo: las instalaciones que afecten al varadero requieren anticipación en una ventana permitida y recepción técnica previa, o una secuencia alternativa que respete H3. El cribado de cortes no resuelve esa dependencia.

\paragraph{Escenario de calendario para desarrollar la red.}
Se utilizará como escenario de análisis un inicio el 1 de diciembre de 2026, sin atribuirlo a una adjudicación o fecha contractual confirmada. H6 cae en diciembre de 2027, H7 en marzo de 2028, H11 en junio de 2028 y H12 en agosto de 2028. Permite estudiar una serie ambiental anterior a 2029, pero su suficiencia requiere medición individual validada, trazabilidad de residuos y evidencia de capacitación ambiental; la fecha por sí sola no cierra las tres no conformidades. La fecha contractual se reconciliará con este escenario antes de aprobar la línea base.

Las restricciones se aplican también a instalación, cambios, pruebas con intervención y capacitación, sin presumir que una actividad está exceptuada por denominarse piloto o trabajo de proveedor. Se programarán las sesiones estacionales antes del congelamiento, con ingreso de los participantes y cohortes acordado con la marina; no se delega capacitación interna en personal TI del cliente. La convivencia que continúa durante el congelamiento no autoriza nuevas intervenciones o sesiones prohibidas. La atención de incidentes conservará su procedimiento contractual, sin convertirlo en autorización para un cambio programado.

\paragraph{Registro de ventanas y reserva.}
Proyecto mantendrá para cada actividad su ámbito, zona, fecha de ejecución, duración, preaviso y responsable de autorización; añadirá los 14 eventos conocidos con un año de anticipación y sus bloqueos locales de uno a tres días. Todo mantenimiento programado requiere aviso al CLIENTE con al menos diez días hábiles, registrado y comprobado antes de autorizar su ejecución. Implantación registrará acceso marítimo a boyas, maniobras protegidas y habilitaciones del cliente. Ante aviso de marejadas se suspenderá inmediatamente toda actividad programada del proyecto, incluso construcción aislada o remota, y se activará la respuesta PL-R01/R-01. Proyecto registrará el efecto sobre todos los frentes, revisará reservas y capacidad y comprobará cese del aviso y ventana admisible antes de reanudar, sin desplazar hitos ni suponer negociable la suspensión \parencites[art.~20, p.~15]{bases-admin}[RT-10.05, p.~22]{bases-transversales}[cap.~10, restricción 12, p.~26; numeral 13.2, p.~29]{caso07}.

La reserva marina de diez días actualmente planteada debe ubicarse en días realmente admisibles anteriores a cada hito. No se cuenta como reserva el congelamiento, una espera contractual ya incluida ni los días ocupados por pruebas o correcciones. El escenario PERT exige además ocho días en E1 y dos en E2 en su subred de datos; esos déficits no desaparecen al cambiar el nombre de la puerta de versión. La selección definitiva requiere una red con duración diaria, disponibilidad y ventanas que permita demostrar ambas reservas, así como recepción y subsanación completas.
